% Appendix A: Scalar, chiral, and Clifford calculations
\label{app:scalar-body}

This appendix records the computations used in \S\ref{sec:scalar-trichotomy} and
\S\ref{sec:six-state}.  Nothing here is new; the purpose is to have the exact
formulas in one place, so that a reader can check the paper's claims without consulting the
source notes.

\subsection{The three quadratic algebras}

For $\kappa\in\{-1,0,1\}$ and $x=a+bj$, $y=c+dj$:
\[
  xy=(ac+\kappa bd)+(ad+bc)j,\qquad
  \bar x=a-bj,\qquad
  N_\kappa(x)=x\bar x=a^{2}-\kappa b^{2}.
\]

\subsubsection{Multiplicativity of the norm}
Expanding both sides,
\[
  N_\kappa(xy)
  =(ac+\kappa bd)^{2}-\kappa(ad+bc)^{2}
  =a^{2}c^{2}+2\kappa abcd+\kappa^{2}b^{2}d^{2}
   -\kappa a^{2}d^{2}-2\kappa abcd-\kappa b^{2}c^{2},
\]
so the cross terms cancel and
\[
  N_\kappa(xy)=a^{2}c^{2}-\kappa a^{2}d^{2}-\kappa b^{2}c^{2}+\kappa^{2}b^{2}d^{2}
  =(a^{2}-\kappa b^{2})(c^{2}-\kappa d^{2})=N_\kappa(x)N_\kappa(y).
\]

\subsubsection{Inverses}
When $N_\kappa(x)=a^{2}-\kappa b^{2}\neq0$,
\[
  x^{-1}=\frac{\bar x}{N_\kappa(x)}=\frac{a-bj}{a^{2}-\kappa b^{2}} .
\]
The excluded loci are therefore:
\[
  \kappa=-1:\ \varnothing\ \text{off the origin};\qquad
  \kappa=0:\ a=0;\qquad
  \kappa=1:\ a=\pm b .
\]

\subsection{The negative-inverse involution}

With $S_\kappa(z)=-1/z$ and $z=a+bj$,
\[
  S_\kappa(a,b)=\Bigl(\frac{-a}{a^{2}-\kappa b^{2}},\ \frac{b}{a^{2}-\kappa b^{2}}\Bigr).
\]
Then
\[
  z\,S_\kappa(z)=(a+bj)\cdot\frac{-a+bj}{a^{2}-\kappa b^{2}}
  =\frac{-a^{2}+abj-abj+\kappa b^{2}}{a^{2}-\kappa b^{2}}=-1 .
\]
For $\kappa=-1$ the denominator is $a^{2}+b^{2}$, so $S_{-1}$ is defined on
$\R^{2}\setminus\{0\}$; and $b>0$ gives $b/(a^{2}+b^{2})>0$, so the upper half-plane is
preserved.  Involutivity is the standard $(S\circ S)(z)=z$ for a M\"obius involution on a
chart where it is defined, checked directly as
\[
  S_{-1}\bigl(S_{-1}(a,b)\bigr)=(a,b).
\]

\subsection{The exchange lift}

With $J=\begin{pmatrix}0&-1\\ 1&0\end{pmatrix}$:
\[
  J^{2}=\begin{pmatrix}-1&0\\ 0&-1\end{pmatrix}=-I,\qquad
  J^{-1}=-J .
\]
Projectively, $[-\id]=[\id]$, hence $[J]^{2}=[I]$ while $J^{2}=-I$ on the lift.

On the chart $z=t/s$, the action $(t,s)\mapsto(-s,t)$ gives
$z\mapsto(-s)/t=-1/z$.

\subsubsection{The generator matrix}
Multiplication by $j$ in the basis $(1,j)$ is
\[
  L_j=\begin{pmatrix}0&\kappa\\ 1&0\end{pmatrix},\qquad
  L_j^{2}=\begin{pmatrix}\kappa&0\\ 0&\kappa\end{pmatrix}=\kappa I,\qquad
  \det L_j=-\kappa .
\]
For $\kappa=-1$: $L_j=J$, $\det=+1$, order four.  For $\kappa=1$: $L_j^{2}=I$,
$\det=-1$.  For $\kappa=0$: $L_j^{2}=0$, $L_j$ nilpotent.

\subsection{Orientation-preservation of the involution}

In real coordinates,
\[
  S(x,y)=\Bigl(\frac{-x}{x^{2}+y^{2}},\ \frac{y}{x^{2}+y^{2}}\Bigr).
\]
As $S$ is the holomorphic map $z\mapsto -1/z$ with derivative $1/z^{2}$,
\[
  \det DS=\left|\frac{1}{z^{2}}\right|^{2}=\frac{1}{(x^{2}+y^{2})^{2}}>0 .
\]

\subsection{The six-state carrier}

Let $T,S,R$ satisfy $T^{2}=S^{2}=R^{2}=1$ and anticommute in pairs, and set
$\Omega=TSR$.

\subsubsection{Centrality}
For example $T\Omega=T\cdot TSR=T^{2}SR=SR$ and
$\Omega T=TSRT=TS(-TR)=-(TST)R=(TTS)R=SR$; the two agree.  The other two generators are
checked identically.

\subsubsection{The square}
\[
  \Omega^{2}=T\,S\,R\,T\,S\,R=T^{2}\,S\,R\,S\,R=S\,R\,S\,R=-S^{2}R^{2}=-1,
\]
where the third equality moves the fourth factor past two factors (sign $+1$) and the
fourth equality moves the third factor past one factor (sign $-1$).

\subsubsection{The pairing}
\[
  \Omega T=SR,\quad \Omega(SR)=-T,\qquad
  \Omega S=RT,\quad \Omega(RT)=-S,\qquad
  \Omega R=TS,\quad \Omega(TS)=-R,
\]
so on each of the three planes $\langle T,SR\rangle$, $\langle S,RT\rangle$, and
$\langle R,TS\rangle$ the operator $J=\Omega\cdot$ acts as
$\begin{pmatrix}0&-1\\ 1&0\end{pmatrix}$.

\subsubsection{The characteristic polynomial}
Since $J$ acts as the same rotation on three complementary two-planes,
\[
  \chi_J(\lambda)=(\lambda^{2}+1)^{3}=\lambda^{6}+3\lambda^{4}+3\lambda^{2}+1 .
\]

\subsubsection{The signature trichotomy}
Allowing $T^{2}=\varepsilon_T$, $S^{2}=\varepsilon_S$, $R^{2}=\varepsilon_R$ multiplies the
above reduction by $\varepsilon_T\varepsilon_S\varepsilon_R$, giving
$\Omega^{2}=-\varepsilon_T\varepsilon_S\varepsilon_R$.  Hence:
\[
  \varepsilon_T\varepsilon_S\varepsilon_R=+1\ \text{(elliptic)},\quad
  -1\ \text{(hyperbolic)},\quad 0\ \text{(parabolic)} .
\]

\subsubsection{The chirality sign law}
A permutation $\pi$ of the generators acts on the word $TSR$ by its sign, so
$P_\pi\Omega P_\pi^{-1}=\sgn(\pi)\Omega$; in particular a transposition gives
$\Omega\mapsto-\Omega=\Omega^{-1}$.

\subsection{The tri-chart rotation}

With
\[
  R=\begin{pmatrix}0&-1\\ 1&-1\end{pmatrix}
\]
one computes $\mathrm{char}_{R}(\lambda)=\lambda^{2}+\lambda+1$, hence $R^{2}+R+I=0$ and
$R^{3}=I$ by Cayley--Hamilton.  In coordinates: $(x,y)\mapsto(-y,x-y)$, and applying this
three times returns $(x,y)$.  The three charts $(x,y)$, $(y-x,-x)$, $(-y,x-y)$ are permuted
cyclically.
